
\subsubsection{5.1 Main Results (RQ1)}

\textbf{Table 1: Pass@1 Results on HumanEval + MBPP (n=664)}

\begin{table}[htbp]
\centering
\resizebox{\columnwidth}{!}{%
\begin{tabular}{ll}
\toprule
Condition & pass@1 \\
\midrule
Static-to-Execution (A) & 55.57\% \\
Execution-to-Static (B) & 43.07\% \\
**Relative Improvement** & **29.02\%** \\
95\% CI & [15.74\%, 44.53\%] \\
McNemar p-value & 6.$31\times 10^{-6}$ \\
\bottomrule
\end{tabular}
}
\end{table}

The static-to-execution ordering achieves a 12.5 percentage point absolute improvement. The 95\% confidence interval excludes zero, and the McNemar test indicates the difference is statistically significant.

\subsubsection{5.2 Mechanism Analysis (RQ2)}

\textbf{Table 2: Early Iteration Gains (H-M1)}

\begin{table}[htbp]
\centering
\resizebox{\columnwidth}{!}{%
\begin{tabular}{ll}
\toprule
Condition & $\Delta Pass_{1}\rightarrow _2$ \\
\midrule
Static-First (A) & 12.50\% \\
Exec-First (B) & 12.05\% \\
Difference & +0.45pp \\
p-value & 0.859 \\
\bottomrule
\end{tabular}
}
\end{table}

Early gains are nearly identical between conditions. The hypothesis that static-first ordering accelerates early gains is not supported.

\textbf{Table 3: Regression Rates (H-M2)}

\begin{table}[htbp]
\centering
\resizebox{\columnwidth}{!}{%
\begin{tabular}{ll}
\toprule
Condition & Regression $Rate_{1}\rightarrow _2$ \\
\midrule
Static-First (A) & 21.53\% \\
Exec-First (B) & 34.69\% \\
Difference & 13.15pp (38\% relative) \\
McNemar statistic & 18.0 \\
p-value & 0.0198 \\
\bottomrule
\end{tabular}
}
\end{table}

Static-first ordering reduces regressions by 38\%. The hypothesis that the ordering effect operates through regression prevention is supported.

\textbf{Per-Iteration Pass Rates:}

\begin{table}[htbp]
\centering
\resizebox{\columnwidth}{!}{%
\begin{tabular}{lll}
\toprule
Iteration & Static-First & Exec-First \\
\midrule
0 & 19.58\% & 17.47\% \\
1 & 51.05\% & 40.81\% \\
2 & 64.61\% & 51.20\% \\
3 & 55.57\% & 43.07\% \\
\bottomrule
\end{tabular}
}
\end{table}

\subsubsection{5.3 Summary}

The ordering effect operates through trajectory stabilization rather than acceleration. Static-first ordering helps the LLM repair more stably by preventing regressions where previously passing tests fail in subsequent iterations.

